\section{7.3 GBDT와 SHAP}
\begin{frame}{다른 조직 방식: \kb{GBDT}}
  \begin{center}
    \small\textbf{``선임이 실수를 남기면, 다음 직원이 바로 그 실수를 고치는'' 순환 근무}
  \end{center}
  \vskip0.4em
  \begin{columns}[T]
    \begin{column}{0.46\textwidth}
      \begin{itemize}
        \item 랜덤 포레스트: 트리를 \textbf{독립적으로} 만들고
          \textbf{다수결}
        \item GBDT: 트리를 \textbf{하나씩 순차적으로} 추가하고,
          \textbf{매번} 새 트리는 ``지금까지의 예측이 아직 못 맞힌 부분
          (\kb{잔차}, residual)''을 목표로 학습
      \end{itemize}
    \end{column}
    \begin{column}{0.5\textwidth}
      \centering
      \includegraphics[width=\linewidth]{ch07_3_rf_vs_gbdt_structure.png}
      \vskip0.3em
      {\footnotesize 랜덤 포레스트(독립·병렬)와 GBDT(순차·잔차 보정)의
      조직 구조 비교.}
    \end{column}
  \end{columns}
\note{\begin{itemize}
\item 7.2의 랜덤 포레스트 = \textbf{``서로 독립적인 전문가 100명의 다수결''.}
  \kb{GBDT}(Gradient Boosted Decision Trees, Friedman 2001)는 전혀 다른
  조직 방식:
\end{itemize}}
\end{frame}
